\section{Selected access and charged tunnels}\label{access:section}

The probabilistic estimate needed later concerns savings selected along
an optimizing path.  A first-moment bound for a fixed local saving does
not control such a selection.  We obtain the stronger estimate by
joining any directed chain of local improvements into one physical path
and comparing its cost with the calibration.

There are three steps.  First, local envelopes control both corrected
access costs and the gauges at cuts, uniformly over the possible cut
occurrences.  Second, a conditional access charge converts a local
negative cost into a tunnel saving that can be realized by a physical
path.  Finally, a sparse corridor network joins every directed chain of
charged tunnels at an expected cost proportional to its span.  This
proves Proposition~\ref{access:collection}, which supplies the collection
hypothesis of Lemma~\ref{probability:thin-tails}.

Throughout, an observation domain for a fixed test is a deterministic
union of independent input cells.  All maxima needed to choose a path,
a cut, or a test are taken before that choice is made.  Conditional
expectations below refer to the fixed tests, and retain local supports.

\subsection{The local input used below}\label{access:local-input}

Fix an admissible pair of input-road rules with common period bound
\[
                       0<\beta\le C_s n^\alpha,
                       \qquad \alpha<1.
\]
The constants $C_s$ and $\alpha$ are fixed as $n$ and the rules vary.
Propositions~\ref{prop:fb-upper} and \ref{prop:fb-positive} place the
near-minimizing roads used in the contraction in this regime.

Road tracking and radius-ten locality are measured in the fixed
$(v,w)$ maximum norm.  We use coordinates in the ordered basis $(A,B)$
below; for $s<1/4$ this change of coordinates has operator norm at most
two.  Comparison with Euclidean coordinates may depend on the fixed flat,
which is harmless here.  Write $\xi=(\xi_1,\xi_2)$ for these coordinates and put
$\ell_1=h(A)>0$, $\ell_2=h(B)>0$.  An $A$-road has nominal equation
$\xi_2=j$, and a $B$-road has nominal equation $\xi_1=i$.  Their tracking
width in these coordinates is at most $\rho=2r<1/100$.

Assign each lattice edge to the half-open unit coordinate square
containing its midpoint.  In each square, collect the assigned edge
weights and independent private seeds into one random element.  The
resulting cell inputs
\[
             (\eta_c:c\in\Z^2)
\]
are independent.  Auxiliary road phases are fixed for the moment; they
are treated explicitly in Section~\ref{access:phases}. Every expectation
and conditional expectation below integrates the spatial cell inputs
with those phase values held fixed. The bounds are uniform over the
fixed admissible phase assignments. Enlarging a support
by a fixed number of unit cells incorporates the dependence neighborhoods
of both road rules and gauges.

We use Lemmas~\ref{found:lem-envelope} and
\ref{found:lem-gauge-tail} for a fixed finite list of road labels near a
cut and bounded parameter windows on those roads.  For each unit cell
$c$, they give nonnegative local envelopes $U_c,V_c$, supported on a
set $\mathcal N(c)$ of at most $m_0$ cells within distance $R_0$ of $c$,
with
\begin{equation}\label{access:envelope-input}
 \E U_c\le C\beta,
 \qquad \E V_c\le a_n,
 \qquad a_n=C_1n^{1-\vartheta},\quad \vartheta>0.
\end{equation}
All constants are uniform in $n$ and the admissible input-road rule.
The two envelopes serve different purposes.  The corrected-cost envelope
$U_c$ bounds the absolute signed cost of travel within the listed windows
and of a change of road label at a physical intersection.  Finite sums
of operations use the corresponding sums of envelopes.  The gauge
envelope $V_c$ bounds the absolute local gauges, after their common
deterministic centering, and the absolute raw signed $Q_n$ costs of the
same short operations.  Both envelopes are functions of the local
inputs.  Their construction requires no local pointwise bound for the
calibration discrepancy $P_n-g$.

The interpolation convention makes these bounds uniform over all
attachment occurrences, including repeated visits to one vertex.  If
$C_j$ is the corrected cost between
successive unit-period anchors, then the corrected signed cost between
parameter values $u,t$ in a bounded interval satisfies
\begin{equation}\label{access:interpolation-envelope}
 |C_W(u,t)|\le
 \sum_{j:\,[j,j+1]\cap [u\wedge t,u\vee t]\ne\varnothing}|C_j|.
\end{equation}
Here each road uses its own parameter, with marked levels at the
integers; the same formula in a shifted wall parameter uses the
correspondingly shifted period endpoints. The number of terms is
bounded by the parameter length plus two. For
switches, the local-gauge lemma supplies the analogous supremum estimate
over the finitely many crossing windows.  Its proof can use calibration
defects, but the envelopes in \eqref{access:envelope-input} themselves are
functions of the local inputs.  Enlarging $U_c,V_c$ by a fixed finite sum
will be understood whenever a cut allows both road families or several
prescribed center labels.

The envelope constants may depend on the fixed flat, the auxiliary row
mesh, and the fixed lists of local labels.  Their support radii are
chosen before the macroscopic piece size $K$ and the number of candidate
shells.  The observation domains used below prove estimates for these
local variables; they are not domains on which the final road rule is
run.

\subsection{Observation domains and conditional charges}

A cut cell $c$ specifies a finite list of road labels and a closed bounded
parameter window on each road.  Endpoints are labeled by path
\emph{occurrences}, so repeated visits to a lattice vertex remain
ordered. For each listed road, prescribe one reference center: its first
occurrence at the absolute nominal longitudinal coordinate of the center
of $c$, using the $\xi_1$ coordinate for an $A$-road and $\xi_2$ for a
$B$-road. This level is fixed by the cell and road label. The
first-occurrence convention is local and commutes with integer cell
translations, and tracking keeps its parameter at bounded distance from
every cut occurrence in $c$. Choose the windows before forming the
envelopes so that they contain these reference centers and all finitely
many allowed row centers. The reference choice adds one center convention
per road label; its half-integer absolute parameter and the fixed finite
wall phase list give only finitely many relative anchor phase classes.

A local test type consists of two cut cells $c_-,c_+$, their road labels
and one prescribed center convention at each endpoint, a reference index
$z\in\Z^2$, and a deterministic closed polygon $D_0$. It must
satisfy
\begin{equation}\label{access:slab}
 \begin{split}
 D_0&\subset [-C K,C K]^2+z,\\
 D_0&\subset
 \{\xi:\ h(c_--z)-b_0\le h(\xi-z)\le h(c_+-z)+b_0\},
 \end{split}
\end{equation}
where a cut cell in an argument of $h$ means its center, and $b_0$ is fixed.
A path of diameter $O(K)$ with backward $h$-excursion at most one fits
in an inner polygon obtained by intersecting a bounding square with the
slab of relative levels $h(c_--z)-b_{\rm in}$ and
$h(c_+-z)+b_{\rm in}$, where the fixed padding $b_{\rm in}$ includes
the unit excursion and the cut-cell diameters. Round the square
coordinates $\xi-z$ and these two relative slab levels outward to a
fixed rational mesh of spacing $\delta_{\rm mesh}>0$, and call the
resulting closed polygon $D_0=z+D_t$. It contains the inner polygon and
satisfies \eqref{access:slab} with
$b_0=b_{\rm in}+\delta_{\rm mesh}$, after a fixed enlargement of $C$.
The rounded lower and upper slab levels remain at bounded distance from
the corresponding cut windows. The bounded relative cut-cell offsets
leave finitely many $D_t$, so these are exact translated types even when
$h(A)$ or $h(B)$ is irrational. Physical access below uses this final
padding $b_0$.

For each relative type, choose a deterministic union $\mathcal D_t$ of
whole input cells before specifying its local phase pattern. Choose it
large enough to determine the weights on $D_t$ and the roads, labels,
and gauges in its cut windows for every admissible local phase pattern
and every input realization. Set $\mathcal D=\mathcal D_{z,t}=z+\mathcal D_t$.
The uniform locality bounds place this stencil in a fixed-width
enlargement of $D_0$. Set
\[
             \mathcal F_{\mathcal D}
                   =\sigma(\eta_c:c\in\mathcal D).
\]
An access envelope may overlap this observation.  The next lemma controls
its conditional mean simultaneously for every possible overlap, without
letting the size of $D_0$ enter the bound.

\begin{lemma}[A local bound for every observation]\label{access:conditioning}
There is a local random variable $W_c$, with the same support as $U_c$,
such that
\begin{equation}\label{access:all-conditionals}
 \E[U_c\mid\mathcal F_{\mathcal D}]\le W_c,
 \qquad \E W_c\le 2^{m_0}C\beta.
\end{equation}
The inequalities hold simultaneously for every deterministic union-cell
observation $\mathcal D$.
\end{lemma}

\begin{proof}
For $I\subseteq\mathcal N(c)$ put
$\mathcal F_I=\sigma(\eta_d:d\in I)$, and define
\[
           W_c=\sum_{I\subseteq\mathcal N(c)}
                        \E[U_c\mid\mathcal F_I].
\]
Independence of the input cells and measurability of $U_c$ give
\[
 \E[U_c\mid\mathcal F_{\mathcal D}]
 =\E[U_c\mid\mathcal F_{\mathcal D\cap\mathcal N(c)}].
\]
This is one of the nonnegative summands defining $W_c$.  There are at most
$2^{m_0}$ summands, each with expectation $\E U_c$.  Since only the
finitely many intersections with $\mathcal N(c)$ matter, versions can be
chosen for which all the inequalities hold simultaneously.
\end{proof}

Only the subsets of $\mathcal N(c)$ enter the construction of $W_c$.
Thus neither its support radius nor its mean depends on the number of
possible $K$-sized observation domains.  We can consequently choose the
separation of candidate cuts before choosing $K$.

A later selection of a test or label uses these simultaneous bounds.
For a chosen deterministic type, its charge is the conditional mean
computed for that type.  We do not condition again on the event that the
type was chosen; this distinction will allow arbitrary selection from
the finite test family.

\subsection{Signed center labels and physical forward access}

Center labels let adjacent pieces use a common additive cost.  The
resulting center adjustments can have either sign.  Before applying
calibration domination, we therefore realize their sum by an actual
path with forward road attachments.

For an oriented road $W$, denote by $C_W(u,t)$ its additive signed
corrected cost, with the gauge of $W$ at both endpoints.  Thus
$C_W(u,t)=-C_W(t,u)$ as bookkeeping, even when the reverse physical trip
has positive cost.  If a physical internal path $\pi$ goes from the
occurrence $x$ on $W_-$ to the occurrence $y$ on $W_+$, define
\[
 C_\pi(x,y)=Q_n(\pi)-g_{W_+}(y)+g_{W_-}(x).
\]
Let $x_0,y_0$ be the chosen center labels on these walls.  The internal
cost measured at center labels is
\begin{equation}\label{access:center-cost}
 H(\pi)=C_{W_-}(x_0,x)+C_\pi(x,y)+C_{W_+}(y,y_0).
\end{equation}

\begin{lemma}[Realization of center-label costs]\label{access:physical}
There are exterior intersection occurrences $a$ on $W_-$ and $b$ on
$W_+$, at bounded parameter distance from their cut windows, such that
$a$ precedes every admissible $x,x_0$, $b$ follows every admissible $y,y_0$,
and the input supports at $a,b$ lie beyond the corresponding observed
$h$-faces.  For every admissible internal path, the physical forward
concatenation $a\to x$, $\pi$, $y\to b$ has corrected cost
\begin{equation}\label{access:physical-identity}
 H(\pi)+C_{W_-}(a,x_0)+C_{W_+}(y_0,b).
\end{equation}
After including changes to exterior network labels, the absolute extra
cost is bounded by $U_{c_-}+U_{c_+}$, enlarging the fixed local envelopes
if necessary.
\end{lemma}

\begin{proof}
The nominal $h$-coordinate increases along either road at a fixed
positive rate, and tracking changes it by a bounded additive amount.
The slab condition therefore lets us pass the full observation halo by
moving a fixed number of periods before the incoming window or after
the outgoing window.  The relevant transverse-road index is chosen from
the deterministic test type.
For example, on an incoming $B$-wall of nominal first coordinate $i$,
take the $A$-road of index
\[
 j_a=\left\lfloor
       \frac{h(c_-)-M_0-\ell_1i}{\ell_2}\right\rfloor,
\]
where the fixed $M_0$ exceeds the final slab padding $b_0$, its input halo,
and the support radius of a connector by a strict margin.  Then use the
canonical intersection of those two prescribed roads.  On an incoming
$A$-wall interchange the coordinates; for exits use the analogous ceiling
above the upper face.  Tracking changes the $h$-level by only $O(\rho)$,
absorbed by the margin.  The chosen parameter differs by a bounded amount
from the cut window.  Thus the index choice uses no unrevealed input,
and all possible attachment occurrences remain in fixed local windows.

The first and last road portions in the asserted concatenation are
physical forward portions.  Additivity of signed corrected costs gives
\[
 \begin{aligned}
 C_{W_-}(a,x)&=C_{W_-}(a,x_0)+C_{W_-}(x_0,x),\\
 C_{W_+}(y,b)&=C_{W_+}(y,y_0)+C_{W_+}(y_0,b).
 \end{aligned}
\]
Adding $C_\pi(x,y)$ proves \eqref{access:physical-identity}, whatever the
relative orders of $x,x_0$ and of $y,y_0$.  In particular, neither signed
center arm has to be traversed backward.  A label change at an exterior
intersection adds the corresponding difference of the two local gauges.
All extra signed travels and switches lie in the bounded windows of
Section~\ref{access:local-input}.  Their absolute costs are therefore
bounded by $U_{c_-}+U_{c_+}$ after a fixed enlargement of those envelopes.
\end{proof}

\subsection{Cuts selected uniformly over optimizing paths}

We next select cuts for which both the corrected access charge and the
raw gauge are small.  The same selection must work for a path chosen
after the whole environment is known.  We achieve this by maximizing
over a finite list of separated candidate tuples.

The union of the two unit-spaced road families divides the plane into
components of uniformly bounded diameter.  Indeed a $B$-road lies in
$|\xi_1-i|\le\rho$ and is proper in its other coordinate.  The planar
crossing property prevents a component of its complement from crossing
this entire strip.  Applying this to the neighboring $B$-roads bounds
the first coordinate of every component; the $A$-roads similarly bound
the second.  Equivalently, inside a large rectangle a left-to-right
crossing intersects every top-to-bottom crossing.  This argument also
applies when roads have self-intersections.  Their nominal parameter
makes the relevant bounded crossing portions finite.

For this cut construction, distances are measured by the maximum norm
in the $(A,B)$ coordinates. Given an initial mark on a parametrized
path, place each successive $K$-exit mark at the first subsequent point
at distance $K$ from the preceding mark, stopping if no such point
remains.

\begin{lemma}[Uniform selected cuts]\label{access:selected-cuts}
Fix a polynomial-box exponent $q$ and an integer $m>4$.  The integer $m$
can be increased so that $m\vartheta>q$. There are fixed radii $H_m$ and
$K>20H_m$, a finite-range field $(E_z)$, and cuts of any
continuous nearest-neighbor path at its successive $K$-exit marks, apart
from bounded end portions, with the following properties.
\begin{enumerate}
\item Each selected cut is a road occurrence within $K/10$ of the
corresponding exit mark.  Consecutive selected cuts have displacement
between $K/2$ and $2K$, and the intervening path piece has diameter $O(K)$.
\item If $c$ is its cut cell and $z$ the unit cell containing its exit
mark, then
\begin{equation}\label{access:cut-control}
             W_c/\beta\le E_z,\qquad V_c\le a_n E_z.
\end{equation}
All finitely many allowed labels and center conventions at that
occurrence are included in these inequalities.
\item Uniformly in the input-road rule and $n$,
\begin{equation}\label{access:cut-tail}
        \PP(E_z>t)\le C_K\min(1,t^{-m}),\qquad t>0.
\end{equation}
In particular $\sup_{n,z}\E E_z^4<\infty$.
\item For any fixed $\varepsilon>0$, the probability that one of these
cuts in a collection of at most $O(n^q)$ possible mark cells has
$V_c>\varepsilon n$ tends to zero.
\end{enumerate}
The construction is valid simultaneously for all paths, including a path
chosen after the entire environment has been observed. At an exit mark,
call the portion after its last preceding hit of the boundary of its
$H_m$-ball the terminal arc. Suppose a fine occurrence partition has
segment traces of diameter at most $\delta_f$ in the present max norm,
and every fine interval meeting that arc chronologically, including at
an endpoint, has displacement in
\begin{equation}\label{access:horizontal-hyp}
 \mathcal G_\kappa=\{z:z_1\ge0,\ |z_2|\le\kappa z_1\},
 \qquad 0\le\kappa<1,
 \qquad H_m>\frac{2(1+\kappa)}{1-\kappa}\delta_f.
\end{equation}
Then all candidates for that mark may be taken on $B$-walls.
\end{lemma}

\begin{proof}
Normalize the two costs together:
\[
                  S_c=W_c/\beta+V_c/a_n.
\]
By \eqref{access:envelope-input} and Lemma~\ref{access:conditioning},
$\E S_c\le C$ uniformly.  Increase the fixed support radius $R_0$
to include the cell diameter.

\emph{Separated candidates.}
Choose $m$ nested annuli about a potential exit mark, with annulus widths
larger than the diameter bound for a component of the road complement,
and with gaps larger than $4R_0$. Also choose an integer
$J>2R_0+1+2\rho$. Take $H_m$ large enough that all annuli fit inside
its ball and $H_m>(m-1)J+2\rho+1$, then take $K>20H_m$.
On the path portion preceding an
exit mark, use the last passage from the outer boundary of this ball to
the mark.  Its successive crossings of the annuli meet the road network;
choose one occurrence at such a meeting in each annulus.  If the mark is
not at a cell center, allow every unit cell meeting the annulus enlarged
by two cell diameters.  The larger gaps ensure disjoint input supports for
any candidate tuple having one cell in each enlarged annulus.  The finite
number of such tuples depends on $K,m,R_0$, but not on $n$.

\emph{One bound for every path.}
For the mark cell $z$, let $\mathcal T_z$ contain every possible annular
tuple and every horizontal tuple described below, allowing all center
offsets within $z$, and set
\[
                  E_z=\max_{(c_1,\ldots,c_m)\in\mathcal T_z}
                                  \min_{1\le j\le m}S_{c_j}.
\]
A finite list of discretized annular and horizontal patterns suffices, because the
allowed candidate sets are subsets of a fixed finite list of unit cells.
For each fixed tuple, independence, Markov's inequality, and the uniform
mean bound give
\[
 \PP(S_{c_1}>t,\ldots,S_{c_m}>t)\le (C/t)^m.
\]
Taking the union over $\mathcal T_z$ proves \eqref{access:cut-tail}.
The actual path supplies a tuple in this list.  Choose the candidate
minimizing $S_c$, breaking ties by the deterministic cell and occurrence
order.  This gives \eqref{access:cut-control} for every path on the same
event, regardless of how that path was selected.

\emph{Geometry and gauges.}
The selected cut lies within $H_m<K/10$ of its mark.  A preliminary exit
has displacement exactly $K$ from the preceding mark, in the continuous
max norm, so the triangle inequality gives the asserted displacement
bounds.  Each portion between preliminary marks stays in its $K$-ball.
The two adjacent short terminal portions give the $O(K)$ diameter bound.
The selected cut occurs after the preceding preliminary mark, since
that mark lies outside the small terminal ball.  Thus the cuts preserve
path order.

Finally,
\[
 \PP(a_nE_z>\varepsilon n)
 \le C_K\left(\frac{a_n}{\varepsilon n}\right)^m
 \le C_{K,\varepsilon}n^{-m\vartheta}.
\]
A union bound over $O(n^q)$ mark cells proves the fourth assertion.
The single objective $S_c$ is essential here: it selects both bounds at
once, so the same product estimate controls the charge and the gauge.

\emph{Horizontal terminal portions.}
Let $[t_-,t_+]$ be the terminal occurrence-time interval and let $P$
linearly interpolate the fine cuts on their time intervals. The diameter
bound gives $|P(t)-\gamma(t)|_\infty\le\delta_f$ in the present
coordinates. For $t_-\le a\le b\le t_+$, the increment $P(b)-P(a)$
is a sum of nonnegative multiples of the good fine increments. Hence,
for $\Delta=\gamma(b)-\gamma(a)$ in these coordinates,
\begin{equation}\label{access:horizontal-submotion}
 \Delta_1\ge-2\delta_f,\qquad
 |\Delta_2|\le\kappa\Delta_1+2(1+\kappa)\delta_f.
\end{equation}
This includes both partial fine segments at the ends. At
$a=t_-,b=t_+$ we have $|\Delta|_\infty=H_m$.
The strict inequality in \eqref{access:horizontal-hyp} excludes
$\Delta_1=-H_m$ and $|\Delta_2|=H_m$, so $\Delta_1=H_m$.
Thus the last boundary hit is on the left face. By lastness and
continuity, all later points of the terminal arc are in the open ball.

Put $p_1=\xi_1(\gamma(t_-))$ and choose the wall levels
\[
 i_j=\lfloor p_1+\rho\rfloor+1+(j-1)J,\qquad 1\le j\le m.
\]
Our choice of $H_m$ gives
$p_1<i_j-\rho$ and $\xi_1(\gamma(t_+))>i_j+\rho$.
A proper $B$-road in each such vertical strip separates these two
sides: apply the planar rectangle-crossing property to a bounded
vertical road portion and a left-to-right subarc. Select the first
meeting along the terminal arc. It is a vertex occurrence, since an
interior meeting of two square-lattice paths lies on a common edge
whose entering vertex would already be a meeting; the initial arc
point is outside the wall. At repeated road visits choose the first
eligible vertex occurrence in the bounded window determined by the
cut cell. The first meetings with $i_1,\ldots,i_m$ occur in that order,
because reaching a later wall requires crossing every earlier wall.
All of them lie strictly inside the terminal ball.

The center of a cell meeting $W_{i_j}$ has first coordinate within
$\rho+1/2$ of $i_j$. Distinct candidate centers are therefore separated
by at least $J-1-2\rho>2R_0$, so their input supports are disjoint.
For a fixed mark cell all candidates lie in a fixed finite
$(H_m+1)$-neighborhood. The horizontal tuples were included in
$\mathcal T_z$ before the common objective was maximized, so the same
product estimate applies without conditioning on the path or on its
terminal test. The selected point stays on the terminal arc, preserving
the preceding order and displacement bounds.
\end{proof}

\subsection{The charged tunnel field}\label{access:tunnels}

Fix one local test type and its observation $\mathcal F_{\mathcal D}$.
The internal improvement is the minimum of the center-label cost over
physical paths in $D_0$ with the specified endpoint windows and labels.
Writing $H(\pi)$ as in \eqref{access:center-cost}, define
\[
                    -u=\beta^{-1}\min_{\pi}H(\pi).
\]
If no admissible path exists, set $u=-\infty$. Otherwise the closed-window
endpoint-selection argument preceding Lemma~\ref{found:lem-envelope}
gives an attained minimum and a measurable local optimizer. It applies
to the finitely many labeled edge occurrences and loop-erased route
types in $D_0$, while retaining the continuum of possible fractions
on each endpoint edge.

The improvement must pay for the two forward attachments of
Lemma~\ref{access:physical}.  Charge their conditional expected cost:
\begin{equation}\label{access:charge}
 a_{\mathcal D}
   =\beta^{-1}\E[U_{c_-}+U_{c_+}\mid\mathcal F_{\mathcal D}],
 \qquad Y=(u-a_{\mathcal D})_+,
\end{equation}
with $Y=0$ when no admissible internal path exists.  Both
$a_{\mathcal D}$ and $Y$ are functions of the observed cells.  Thus $Y$
is a local saving after access has been paid for.  Its sign comes from
the additive center-label cost; its realization will use physical
forward attachments.

More explicitly, for a fixed local phase pattern $\phi$, the charge is
the kernel
\[
 a_{\mathcal D}^{\phi}(\eta)
 =\beta^{-1}\E_{\eta}
   [U_{c_-}^{\phi}+U_{c_+}^{\phi}\mid\eta_{\mathcal D}],
\]
where $\E_\eta$ integrates spatial inputs only. We suppress $\phi$
while it is fixed. For the original random phase field $\Phi$, evaluate
these same kernels and set
$Y(\eta,\Phi)=(u^{\Phi}(\eta)-a_{\mathcal D}^{\Phi}(\eta))_+$.
For each relative type and phase pattern, define a product-integral
version of this kernel at the reference site by integrating the
unobserved spatial coordinates against their product law. Define its
copies by translation with $z$, and use these same kernels for the
original and template fields. This fixes the versions covariantly.

The physical-competitor envelope of Lemma~\ref{found:lem-envelope},
with this domain and these center labels, gives
$Y\le u_+\le U'/\beta$ and hence $\E Y\le C$. In particular, the
maximum reward over any finite array is integrable, being bounded by
the sum of its finitely many $Y$ variables.

\begin{lemma}[Tunnel accounting]\label{access:accounting}
If a path piece with the specified type has normalized center-label cost
$c$, then
\begin{equation}\label{access:negative-accounting}
            c^-\le Y+W_{c_-}/\beta+W_{c_+}/\beta.
\end{equation}
If $Y>0$, the internal minimizer has a physical forward realization
between its exterior intersection labels whose conditional expected
normalized corrected cost is at most $-Y$.
\end{lemma}

\begin{proof}
Minimality gives $c\ge-u$, and $a_{\mathcal D}\ge0$ gives
\[
 c^-\le u_+\le (u-a_{\mathcal D})_++a_{\mathcal D}.
\]
Lemma~\ref{access:conditioning} bounds $a_{\mathcal D}$ by the two
$W/\beta$ terms in \eqref{access:negative-accounting}.

On $\{Y>0\}$, Lemma~\ref{access:physical} realizes the internal
minimizer with normalized corrected cost at most
$-u+(U_{c_-}+U_{c_+})/\beta$.  Taking its conditional mean yields
$-u+a_{\mathcal D}=-Y$.  The envelopes already bound every admissible
attachment occurrence, so this estimate holds for the selected minimizer.
\end{proof}

At selected cuts, \eqref{access:negative-accounting} reads
\[
                    c^-\le Y+E_{z_-}+E_{z_+}.
\]
The error has the light tail \eqref{access:cut-tail}.  To assign one error
to the piece center, sum the finitely many $E_z$ in its $O(K)$
neighborhood; this preserves locality and the tail exponent.

The remaining task is to control selected tunnel rewards $Y$.  A piece
may choose its type after inspection: we bound the resulting sums by
summing the estimates for the finite type list.  This needs no
independence between types at one site.

\subsection{Compatible deterministic phase templates}\label{access:phases}

To use the abstract thin-diagram estimate, we need an iid copy of each
local tunnel law on a sparse lattice.  Shared wall phases are fixed
first; we will then dominate the original random-phase field by a finite
sum of such template fields.

Fix the rational mesh, the road periods, and the finitely many row and
anchor phase classes used by a local test.  Choose an even integer $S$ that is
a multiple of every resulting period and is much larger than the diameter
of $\mathcal D$ plus every dependence radius.  In $(A,B)$ coordinates
place copies of the test at
\[
                        p_{ij}=S(i,j),\qquad (i,j)\in\Z^2.
\]
If two translated copies refer to the same infinite $A$-road, their
longitudinal translations differ by an integer multiple of $S$; their
phase prescriptions therefore agree.  The same is true of shared
$B$-roads.  The finitely many transverse offset classes used by the test
can be repeated periodically.  Thus one fixed local phase template has a
consistent extension to the whole road network.  Give all remaining
roads fixed phases with the same periodicity.

Once the phases are deterministic, the copied observation sets have
disjoint cell supports, including all private seeds.  Translation
covariance gives a common law, and independence of the cells makes
$(Y_{ij})$ iid.  Shared phases along one wall have introduced no additional
randomness into this assertion.

For the original construction, enumerate the finitely many local phase
templates, together with the needed coarse residue and row classes, and
apply the collection estimate to each.  The local moment bounds hold
uniformly for every fixed admissible phase assignment.  The enumeration
and $S$ may depend on the fixed mesh, type, and road periods, but are
fixed before $n$ varies.

The domination is pointwise in one common coupling.  A template records
every wall phase used by the test and its access envelopes, as well as
the phases of both anchors in every projection in
\eqref{found:eq-gauge-formula}.  This includes the reference anchors
$z_S$.  Lemma~\ref{phases:locality} bounds the resulting finite list
uniformly.
For each sparse residue class and each assignment $t$ to the list,
repeat $t$ periodically and evaluate the field $Y_z^t$ using the
original weights and the same road-indexed spatial private inputs.
Only the phase arguments are overridden. At each site $z$, one
assignment is its original local phase pattern; locality then gives
$Y_z=Y_z^t$ for that assignment, and hence
\[
                         Y_z\le\sum_tY_z^t.
\]
For fixed $t$ the sparse field is iid; the different template fields
need not be independent. The same coupling applies to local error
fields. The conditional expectations defining gauges and access charges
integrate weights and private coordinates only. Thus the full local phase
list determines their kernels, including projected terms and anchor
means, and the matching template uses exactly the original charge.

\subsection{Corridors avoiding every observed halo}

We now build the routes used in the collection estimate.  Their ordinary
connectors must avoid every observed input, including observations at
sites that the chosen chain does not use.  This ensures that choosing the
chain does not bias the expected connector costs.

Increase $S$ by an integer factor so that the observation at $p$ and the
supports of its two charged access envelopes lie in
\[
                     p+[-S/10,S/10]^2.
\]
The physical trace and full input support of each ordinary road segment,
and the full input support of each intersection cost, lie within distance
$R_*$ of the corresponding nominal segment or intersection.
Choose
$S>100R_*$.  Reserve the $A$-roads on
$\xi_2=(j+\tfrac12)S$ and the $B$-roads on
$\xi_1=(i+\tfrac12)S$.  Their entire input-support strips miss every
observation square.  Adding these sparse road copies does not change
any already specified road.  At each crossing choose a canonical
intersection occurrence.  Since tracking is much smaller than the
spacing, the intersections with consecutive transverse roads occur in
the forward order.  A monotone lattice path on this corridor grid is
therefore an actual forward concatenation of roads.

The southwest and northeast portals of the cell centered at $p_{ij}$
are the intersections near
\[
       p_{ij}-(S/2,S/2),\qquad p_{ij}+(S/2,S/2),
\]
respectively.  Connect its southwest portal to the exterior entry $a$
from Lemma~\ref{access:physical} as follows.  If the last incoming road is
a $B$-wall, approach its transverse $A$-road along the left corridor
boundary and then along that $A$-road to $a$.  If the last incoming road
is an $A$-wall, interchange the two coordinates and use the bottom
boundary.  The final transverse portion is below the lower observed
$h$-face: $a$ was chosen beyond that face by more than the dependence
radius, and every earlier point of a forward $A$- or $B$-segment has
no larger nominal $h$, including during pauses. All preceding boundary portions are outside the
observation square.  The analogous construction from $b$ to the northeast
portal uses the upper observed $h$-face and the top or right boundary.

The nominal segments of the ordinary connector portions assigned to this
cell stay in $p+[-S/2,S/2]^2$ and have total length $O(S)$. Their physical
traces and full input supports lie in its $R_*$-enlargement. Apart from the two short charged
attachments in \eqref{access:charge}, their input supports avoid the
cell's observation. They also avoid every other observation: the
$R_*$-enlargement misses the inner observation squares because
$S>100R_*$. Shift a corridor or extend a short access
by a fixed extra period whenever necessary to preserve a full input
buffer.  The available margins are fixed fractions of $S$, so these
changes preserve the forward ordering and the $O(S)$ length bound.

Figure~\ref{fig:observed-halos-corridors} separates the ordinary route
portions, whose input supports avoid every observation, from the two
charged attachments at each selected cell.

\begin{figure}[t]
\centering
\begin{tikzpicture}[x=1cm,y=1cm,>=Latex,
  every node/.style={font=\scriptsize},
  grid road/.style={draw=black!28,line width=.35pt},
  ordinary/.style={draw=black,line width=.8pt},
  charged/.style={draw=black,line width=1pt,densely dashed},
  tunnel/.style={draw=black,line width=1.35pt},
  halo/.style={draw=black!48,fill=black!12,dashed,line width=.4pt},
  portal/.style={circle,draw=black,fill=white,inner sep=1.25pt}]

\begin{scope}[shift={(.35,.8)}]
  \node[font=\footnotesize] at (2.7,5.83)
    {(a) A route through two selected cells};
  \foreach \v in {0,1.8,3.6,5.4} {
    \draw[grid road] (\v,0)--(\v,5.4);
    \draw[grid road] (0,\v)--(5.4,\v);
  }
  \foreach \sx/\sy in {0/0,1.8/0,3.6/0,0/1.8,1.8/1.8,3.6/1.8,
                         0/3.6,1.8/3.6,3.6/3.6} {
    \begin{scope}[shift={(\sx,\sy)}]
      \path[halo] (.67,.91)--(.91,.67)--(1.13,.67)--(1.13,.89)
        --(.89,1.13)--(.67,1.13)--cycle;
    \end{scope}
  }
  \foreach \sx/\sy in {0/0,3.6/3.6} {
    \begin{scope}[shift={(\sx,\sy)}]
      \draw[ordinary,-{Latex[length=1.2mm]}]
        (0,0)--(0,.75)--(.76,.75);
      \draw[charged,-{Latex[length=1.2mm]}] (.76,.75)--(.76,.88);
      \draw[tunnel,-{Latex[length=1.2mm]}]
        (.76,.88)--(.84,.91)--(.90,.84)--(.96,.94)--(1.04,.93);
      \draw[charged,-{Latex[length=1.2mm]}] (1.04,.93)--(1.04,1.05);
      \draw[ordinary,-{Latex[length=1.2mm]}]
        (1.04,1.05)--(1.8,1.05)--(1.8,1.8);
      \node[portal] at (0,0) {};
      \node[portal] at (1.8,1.8) {};
    \end{scope}
  }
  \draw[ordinary,-{Latex[length=1.4mm]}]
    (1.8,1.8)--(3.6,1.8)--(3.6,3.6);
  \draw[ordinary,-{Latex[length=1.4mm]}] (-.3,0)--(0,0);
  \draw[ordinary,-{Latex[length=1.4mm]}] (5.4,5.4)--(5.7,5.4);
  \node[anchor=north east] at (-.02,-.04) {SW};
  \node[anchor=south east] at (1.76,1.86) {NE};
  \node[anchor=north] at (.9,.47) {selected};
  \node[anchor=south] at (4.5,4.83) {selected};
  \node[align=center] at (2.7,3.23) {unused\\ observation};
  \draw[black!55,line width=.35pt] (2.7,3.01)--(2.7,2.9);
  \draw[<->,black!65,line width=.35pt] (0,-.36)--(1.8,-.36);
  \node[fill=white,inner sep=1pt] at (.9,-.36) {$S$};
\end{scope}

\begin{scope}[shift={(7.15,.8)}]
  \node[font=\footnotesize] at (2.5,5.83)
    {(b) Physical access in one selected cell};
  \draw[grid road] (0,0) rectangle (5,5);
  \path[halo] (1.3,2.05)--(2.05,1.3)--(3.6,1.3)--(3.6,2.7)
    --(2.8,3.5)--(1.3,3.5)--cycle;
  \draw[black!48,line width=.45pt]
    (1.48,2.12)--(2.12,1.48)--(3.42,1.48)--(3.42,2.62)
    --(2.72,3.32)--(1.48,3.32)--cycle;
  \node[black!70] at (2.65,1.68) {$D_0$};
  \node[black!70] at (2.18,3.05) {observed halo};
  \draw[black!35,line width=.45pt] (1.7,1.06)--(1.7,2.16);
  \draw[black!35,line width=.45pt] (3.15,2.56)--(3.15,3.73);
  \draw[ordinary,-{Latex[length=1.6mm]}]
    (0,0)--(0,1.4)--(1.7,1.4);
  \draw[charged,-{Latex[length=1.6mm]}] (1.7,1.4)--(1.7,1.95);
  \draw[tunnel,-{Latex[length=1.6mm]}]
    (1.7,1.95)--(1.93,2.05)--(2.12,1.93)--(2.36,2.32)
    --(2.65,2.2)--(2.85,2.62)--(3.15,2.78);
  \draw[charged,-{Latex[length=1.6mm]}] (3.15,2.78)--(3.15,3.4);
  \draw[ordinary,-{Latex[length=1.6mm]}]
    (3.15,3.4)--(5,3.4)--(5,5);
  \node[portal] at (0,0) {};
  \node[portal] at (5,5) {};
  \foreach \px/\py in {1.7/1.4,1.7/1.95,3.15/2.78,3.15/3.4}
    \fill (\px,\py) circle (1.15pt);
  \node[anchor=north west] at (.04,-.08) {southwest portal};
  \node[anchor=south east] at (4.96,5.04) {northeast portal};
  \node[anchor=north] at (1.7,1.35) {$a$};
  \node[anchor=west] at (1.8,1.84) {$x$};
  \node[anchor=west] at (3.23,2.75) {$y$};
  \node[anchor=south] at (3.15,3.46) {$b$};
  \node[anchor=north] at (2.47,2.12) {$\pi$};
  \node[anchor=east] at (1.59,1.65) {$W_-$};
  \node[anchor=west] at (3.27,3.1) {$W_+$};
  \node[anchor=east,align=right] at (1.05,2.22)
    {lower\\ $h$-face};
  \draw[black!55,-{Latex[length=1mm]},line width=.35pt]
    (1.11,2.14)--(1.44,1.91);
  \node[anchor=west,align=left] at (3.65,4.03)
    {upper\\ $h$-face};
  \draw[black!55,-{Latex[length=1mm]},line width=.35pt]
    (3.65,3.89)--(3.06,3.24);
  \draw[-{Latex[length=1.4mm]},black!60,line width=.45pt]
    (4.15,.55)--(4.72,1.12);
  \node[black!70,anchor=north] at (4.43,.48) {$h$ increases};
\end{scope}

\draw[ordinary] (.45,-.05)--(1.05,-.05);
\node[anchor=west] at (1.12,-.05) {ordinary connector};
\draw[charged] (4,-.05)--(4.6,-.05);
\node[anchor=west] at (4.67,-.05) {charged attachment};
\draw[tunnel] (8,-.05)--(8.17,.01)--(8.32,-.11)--(8.6,-.05);
\node[anchor=west] at (8.67,-.05) {internal tunnel $\pi$};
\end{tikzpicture}
\caption{Physical routing through selected cells, schematically and not to
scale. Gray regions represent observed input halos, including observations
at unused cells. The ordinary connectors have input supports disjoint from
every observed halo. The two dashed attachments at a selected tunnel are
also physical forward portions; their supports may meet only that cell's
observation. In the enlarged cell, the inner outline is $D_0$, the exterior
entry $a$ is below the lower
observed $h$-face and the exterior exit $b$ is above the upper face. The
transverse connectors can therefore reach them without meeting the
observation. The internal path is $\pi:x\to y$. Only physical route portions
are drawn: the signed center-label terms in $H(\pi)$ are bookkeeping and do
not prescribe additional traversals.}
\label{fig:observed-halos-corridors}
\end{figure}

The construction works for every allowed incoming and outgoing label.
After conditioning on all observations, the selected labels and nominal
itineraries are fixed, whereas their actual road geometry and
intersection occurrences remain random.  For each itinerary, use the
envelopes over its full prescribed crossing windows.  Its uncharged
support avoids all observations, so those envelopes retain their
unconditional means.  The mean bound is uniform over the finite label
list and therefore also applies to the observation-dependent choice.

\begin{proposition}[Uniform finite-volume collection bound]
\label{access:collection}
For the iid tunnel field associated with any fixed local type and phase
template there are constants $C,C_0<\infty$, independent of $n$ and the
input-road rule, such that the following holds.  Let $G_n(L)$ be the
maximum of $\sum_{p\in\Gamma}Y_p$ over finite distinct chains in an
arbitrary deterministic translated coordinate square of side at most
$L$, whose consecutive differences have both coordinates at least $2S$.
The empty chain and singletons are allowed.  Then
\begin{equation}\label{access:finite-volume}
                       \E G_n(L)\le CL+C_0.
\end{equation}
The same bound therefore holds if consecutive differences are additionally
required to lie in any fixed inner cone of the positive quadrant.
\end{proposition}

\begin{proof}
\emph{Choose the chain and realize its route.}
Observe the inputs of every test in the sparse array.  In the given
finite square, choose a reward-maximizing chain, resolve ties
lexicographically, and omit zero-reward sites.  This is a function of the
observations.  Fix an
external southwest and an external northeast corridor intersection,
outside the square by $3S$ in both coordinates, with the same prescribed
intersection-label convention.  These endpoints depend on the square,
not on the observed chain.

Route monotonically on the corridor grid from the start to the southwest
portal of the first selected cell, use that tunnel and its private access,
then continue from its northeast portal to the next southwest portal.
The assumed $2S$ gaps make both coordinates of each latter displacement
nonnegative, with a full corridor spacing to spare.  Continue similarly
to the fixed end.  Every portion is a physical forward road portion or a
physical internal tunnel path.  Thus the result is an actual path.

Its total ordinary-corridor length is $O(L+S)$, by telescoping the two
positive coordinate increments.  A chain has at most $C(L/S+1)$ sites.
The private uncharged connector lengths are $O(S)$ per site, so their
total length is also $O(L+S)$.  The number of road switches has the same
bound when measured in unit road spacings.  The two short charged
attachments at each site are excluded from this count.

All selected route costs are integrable. A used tunnel has $Y>0$ and
therefore has negative center-label cost. Rearranging that cost identity
bounds its positive physical passage cost by a bounded $nh$ displacement
and the integrable local gauge and raw-arm envelopes. Charged attachments
and ordinary finite road portions are integrable by the same window
estimates. The chain is finite, and its maximum reward is integrable by
the bound following \eqref{access:charge}. Thus the conditional costs
and the expectation rearrangements below are finite.

\emph{Take conditional expected costs.}
For every prescribed ordinary connector, its full input support misses
\emph{all} observations.  Its inputs therefore retain their product law
under the conditioning.  Subdivide the connector into unit road periods
and apply the local envelopes to those periods and its switches.  Its
expected normalized corrected cost is at most a constant times its
nominal length plus its number of switches.  Summing uses only linearity
of conditional expectation; connectors may share unobserved inputs.

The support of a charged access envelope meets no observation other
than its own test.  Independence gives
\[
 \E[U_{c_-}+U_{c_+}\mid\hbox{all observations}]
 =\E[U_{c_-}+U_{c_+}\mid\mathcal F_{\mathcal D_p}].
\]
Thus each selected tunnel together with its charged attachments has
conditional expected normalized corrected cost at most $-Y_p$, by
Lemma~\ref{access:accounting}.  Consequently the completed path $\Pi$
satisfies
\begin{equation}\label{access:conditional-route}
 \E\left[\frac{Q_n(\Pi)-\Delta g(\Pi)}{\beta}
           \,\middle|\,\hbox{all observations}\right]
                 \le C(L+S)-G_n(L).
\end{equation}

\emph{Apply the calibration at the fixed endpoint windows.}
Corrected costs use only differences of gauges. For this endpoint
calculation, undo any common deterministic centering and use the
original mean-matching gauge representatives; the corrected cost of
$\Pi$ is unchanged.
The constructed path has endpoint labels prescribed by translated copies
of the exterior anchor convention.  Calibration domination gives,
pathwise,
\[
 Q_n(\Pi)-\Delta g(\Pi)
 \ge (P_n-g)(x_{\mathrm{end}})-(P_n-g)(x_{\mathrm{start}}).
\]
Although an endpoint intersection is random, it lies in a prescribed
bounded road window.  Put $R_W=P_n-g_W$ on that exterior road and choose
a fixed marked anchor $k$ before the entire window of possible endpoint
occurrences $t$.  Equation~\eqref{found:eq-mean} gives
$\E R_W(k)=0$, and the defect and interpolation identities give
\[
 R_W(t)-R_W(k)=-D_W(k,t)+C_W(k,t).
\]
The interval is covered by a fixed number of complete road periods.
Additivity and nonnegativity bound its defect by the defects of those
periods; Lemma~\ref{found:lem-envelope} controls its corrected cost.
Thus the expected supremum of the absolute right side is at most
$C\beta$.  In particular,
\[
 \bigl|\E(P_n-g)(x_{\mathrm{start}})\bigr|
 +\bigl|\E(P_n-g)(x_{\mathrm{end}})\bigr|
 \le C\beta.
\]
The estimate covers every possible intersection occurrence and allowed
endpoint label.  Consequently the expected normalized corrected cost of
$\Pi$ is at least $-C$.  Taking expectations in
\eqref{access:conditional-route} yields
$\E G_n(L)\le C(L+S)+C$, and the fixed $S$ is absorbed into $C_0$.
The bound holds at each finite $L$, uniformly over the tunnel law; no
law-dependent limiting passage is needed.
\end{proof}

The constants in Proposition~\ref{access:collection} may depend on the
fixed type, mesh, and sparse spacing, and are uniform over all $n$ and
road rules satisfying the local input assumptions.  Taking their maximum
over the finite type list gives one bound for that list.  For rational
$v,w$ in fixed sufficiently small neighborhoods of the flat endpoints,
the coordinate comparisons and the positive rates $h(A),h(B)$ have
uniform bounds.  The domains, support buffers, and sparse spacing can
therefore be fixed throughout those neighborhoods.  In the application
in Section~\ref{links:section}, this test family is fixed before the
final thin-diagram cone is narrowed.  Rational direction denominators
affect only the divisibility required of $n$.

Together with the phase-template domination, the proposition supplies
the iid collection bounds required by Lemma~\ref{probability:thin-tails}.
The observation domains, candidate shells, and corridors enter this
distributional proof; they do not enlarge the input support of the final
road-improvement rule.
